\documentclass[aspectratio=169, 11pt]{beamer}
\usepackage{beamer_preamble}

\begin{document}

% ==== Title ===============================================================
\titleslide{Lesson 3-5 \textperiodcentered\ Period 2}{Multiplicity and Factored-Form Graphing}

% ==== Do Now ==============================================================
\begin{frame}{Do Now --- Find the zeros from factored form}{Algebra 2 \textperiodcentered\ Lesson 3-5 \textperiodcentered\ Period 2}
\phasebadges{1}{5}{bridge from Monday}
\vspace{0.2cm}
\begin{projcard}{Practice (3-5-savvas-q13) \textperiodcentered\ Sketch from zeros}{slidegold}
Sketch the graph of the function by finding the zeros:
\[g(x) = (x+3)(x-1)(x-4)\]

\vspace{0.1cm}
For each factor: what zero does it give? Mark each zero on your sketch.\\
Use the sentence frame: ``The zeros are \(x=\) \underline{\hspace{0.6cm}}, \(x=\) \underline{\hspace{0.6cm}},
and \(x=\) \underline{\hspace{0.6cm}} because each factor equals zero.''
\end{projcard}
\end{frame}

% ==== Launch ==============================================================
\begin{frame}{Launch --- Multiplicity Rules (Example 2)}{Algebra 2 \textperiodcentered\ Lesson 3-5 \textperiodcentered\ Period 2}
\phasebadges{2}{12}{teacher models}
\vspace{0.2cm}
\begin{projcard}{Multiplicity Rules --- keep on your page all period}{slidegreen}
\begin{enumerate}
  \setlength{\itemsep}{0pt}
  \item If a factor appears \textbf{ODD} number of times \(\rightarrow\) the graph \textbf{CROSSES} the x-axis at that zero.
  \item If a factor appears \textbf{EVEN} number of times \(\rightarrow\) the graph is \textbf{TANGENT} (touches and turns back) at that zero.
\end{enumerate}

\vspace{0.1cm}
\small
Example (Savvas Ex 2): \(f(x) = (x-2)^3(x+1)^2\)\\
\hspace*{1em}\(\bullet\) \(x = 2\): factor appears 3 times --- ODD \(\rightarrow\) \textbf{crosses}\\
\hspace*{1em}\(\bullet\) \(x = -1\): factor appears 2 times --- EVEN \(\rightarrow\) \textbf{tangent}

\vspace{0.1cm}
\textcolor{slideaccent}{\textbf{Count the factor. Odd or even? That's your answer.}}
\end{projcard}
\end{frame}

% ==== Explore =============================================================
\begin{frame}[allowframebreaks]{Explore --- TPS (Think-Pair-Share)}{Algebra 2 \textperiodcentered\ Lesson 3-5 \textperiodcentered\ Period 2}
\phasebadges{2}{33}{student work}
\small\textcolor{slidegray}{5 items in order. For each zero: write multiplicity + crosses/tangent before sketching. \(\sim\)6--7 min each.}

\vspace{0.1cm}
\begin{projcard}{Try It 2a \textperiodcentered\ Describe zeros}{slideblue}
\(f(x) = x(x+4)(x-1)^4\) --- describe graph behavior at each zero.
\end{projcard}
\vspace{-0.15cm}
\begin{projcard}{Try It 2b \textperiodcentered\ Real zeros only}{slideblue}
\(f(x) = (x^2+9)(x-1)^5(x+2)^2\) --- which zeros are real? Crosses or tangent?
\end{projcard}
\vspace{-0.15cm}
\begin{projcard}{Practice \#14 \textperiodcentered\ Factor first}{slideblue}
\(f(x) = x^3 - 8x^2 + 16x\) --- find zeros, then describe graph behavior.
\end{projcard}
\vspace{-0.15cm}
\begin{projcard}{Practice \#15 \textperiodcentered\ Factor by grouping}{slideblue}
\(g(x) = x^3 - x^2 - 25x + 25\) --- find zeros, then describe graph behavior.
\end{projcard}
\vspace{-0.15cm}
\begin{projcard}{Practice \#16 \textperiodcentered\ Biquadratic}{slideblue}
\(f(x) = 9x^4 - 40x^2 + 16\) --- let \(u = x^2\); find all zeros.
\end{projcard}
\end{frame}

% ==== Share / Summary =====================================================
\begin{frame}{Share / Summary}{Algebra 2 \textperiodcentered\ Lesson 3-5 \textperiodcentered\ Period 2}
\phasebadges{2}{5}{academic conversation + bridge prompt}
\vspace{0.2cm}
\begin{projcard}{Sentence frame \textperiodcentered\ Pick one of your zeros}{slideblue}
``For \(f(x) = \) \underline{\hspace{2cm}}, the zero \(x = \) \underline{\hspace{1cm}}
has multiplicity \underline{\hspace{0.8cm}},\\
so the graph \underline{\hspace{2.2cm}} the x-axis at that zero.''
\end{projcard}

\vspace{0.15cm}
\begin{projcard}{Bridge to Wednesday (Period 3 Do Now)}{slideaccent}
A cubic has zeros \(2\), \(1+\sqrt{3}\), and \(1-\sqrt{3}\).\\
What is the nature of its coefficients --- rational, irrational, or complex?\\
\textcolor{slideaccent}{\textbf{Write the key term ``conjugate pair'' in your packet margin.}}\\
\small(Directly prepares Topic 3 Assessment Q\#2.)
\end{projcard}
\end{frame}

% ==== Exit Ticket =========================================================
\begin{frame}{Exit Ticket --- Summary Recap}{Algebra 2 \textperiodcentered\ Lesson 3-5 \textperiodcentered\ Period 2}
\phasebadges{1-2}{2}{not graded}
\vspace{0.2cm}
\begin{projcard}{Complete on your exit ticket}{slidegold}
\begin{enumerate}
  \setlength{\itemsep}{0.2cm}
  \item Today I learned that \underline{\hspace{6cm}}.
  \item The rule I want to remember is \underline{\hspace{5cm}}.
  \item One thing I'm still unsure about is \underline{\hspace{4.5cm}}.
\end{enumerate}
\end{projcard}
\end{frame}

\end{document}
